\subsection{Root Mean Squared Error}
É um estimador da diferença entre o valor observado e o valor real; Dado pela seguinte fórmula $$ RMSE =
\sqrt{\frac{1}{n}\sum_{i=1}^n(Y_i - Y_i')^2}$$ Yi são os valores reais e Y'i são os valores observados. Além disso, *n*
é o número de amostras. Aquilo que é medido na verdade é a média do erro ao quadrado, ou seja, outliers são punidos mais
fortemente do que outros pontos, pois quanto mais longe estiver, maior será o RMSE;


\subsection{Mean Absolute Error}


\subsection{Nash-Sutcliffe Model Efficiency Coefficient}
Nash Sutcliffe efficiency -> é uma Machine Learning Metrics|Métrica que compara a magnitude residual do resíduo em
comparação com a variação dos dados medidos. Ou seja, ele server para compara dados simulados (modelo) com dado medido
(real). É dado por (Com $Q_0$ como média dos dados observados e $Q_m$ os dados modelados e $Q^t_0$ os dados observados
em um dado tempo):

\[
\begin{aligned}
NSE = 1 - \frac{\sum^T_{t=1} (Q^t_0 - Q^t_m)^2}{\sum^T_{t=1} (Q^t_0 - Q_0)^2}
\end{aligned}
\]

Essa métrica é normalizada entre 0 e 1. Para o caso de 1, o modelo perfeitamente descreve os dados observados. Contudo,
o que é interessante dessa métrica é para modelos com NSE 0. Para esses casos, o modelo teve a mesma habilidade
preditiva do que a média de uma Série Temporal em relação à soma do erro quadrado; * Valores mais próximos de 1 sugerem
uma habilidade preditiva maior; Como sempre, essa normalização não é perfeita. Para casos com NSE<0, a média do
observado é um melhor preditor do que o modelo;
